\RequirePackage[utf8]{inputenc}
\documentclass[soumission]{jsfds}
% \usepackage[latin1]{inputenc}
%\usepackage[utf8]{inputenc}

\usepackage[T1]{fontenc}
\usepackage{hyperref} 
\usepackage{amsmath,amsthm} 
\usepackage{graphicx}
\usepackage{natbib} 

\startlocaldefs
% Commandes auteurs spécifiques, par exemple :
\newcommand{\R}{\mathbb{R}}
\newcommand{\E}{\mbox{\tiny \`eme}}
\def\1{1\kern-.20em {\rm l}}
\endlocaldefs

\setmainlanguagefrench
\firstpage{1} % à régler par l'éditeur

\begin{document}

\begin{frontmatter}

  % infos fournies par l'éditeur, laisser vide
  % \volume{}
  % \numero{}
  % \anneepublication{}
  % \doi{}
 \arxiv{math.PR/00000000} % Si disponible

  \titre{Manuel d'utilisation de la classe de document du J-SFdS\thanks{Merci à S\'ebastien Mengin (edilibre.net) pour l'adaptation de la classe de documents \texttt{imsart}.}}
  \titrecourant{Manuel du J-SFdS}
  \titrealter{Authors users guide of the J-SFdS \LaTeX document class}

  \begin{InfosAuteurs}
    \auteur{%
      \prenom{Auteur} 
      \nom{Un}%
      \thanksref{t1}%
      \contact[label=e1]{auteur.un@math.univ.fr}%
    }%
    \and%
    \auteur{%
      \prenom{Auteur} 
      \nom{Deux}%
      \thanksref{t2}%
      \contact[label=e2]{auteur.deux@stat.edu}%
    }%
    \and%
    \auteur{%
      \prenom{Auteur} 
      \nom{Trois}%
      \thanksref{t2}%
      \contact[label=e3]{auteur.trois@stat.edu}%
    }

    \affiliation[t1]{Nom de l'institution auteur 1.\\ 
      \printcontact{e1}%
    }

    \affiliation[t2]{Nom de l'institution des auteurs 2 et 3.\\ 
      \printcontact{e2} and \printcontact{e3}% 
    }

    \enteteauteurs{Un, Deux et Trois} %noms courants des auteurs en tête de page
  \end{InfosAuteurs}

  
  \begin{abstract}
Cette courte note décrit la classe \LaTeX \texttt{jsfds} et illustre son usage en se présentant sous la
forme d'un article du Journal de la Société Française de Statistique. Cette classe est une adaptation de la classe publique \texttt{imsart}. Le résumé du contenu de l'article doit être clair, decriptif, auto-suffisant et pas plus long que 200 mots. Il doit aussi être adapté à publication dans les recueil de résumés. Il est préférable d'éviter des formules mathématiques.
\end{abstract}

  
  \begin{altabstract}
This paper describes the use of the \texttt{jsfds} \LaTeX document class and is prepared as a sample
to illustrate the use of this class written for the Journal of the French Statistical Society. This an adaptation of the public document class \texttt{ imsart}. The abstract should summarize the contents of the paper. It should be clear, descriptive, self-explanatory and not longer than 200
words. It should also be suitable for publication in abstracting services.
Please avoid using math formulas as much as possible.
\end{altabstract}


  \begin{motscles}
    \mot{mode d'emploi}%
    \mot{classe du J-SFdS}%
  \end{motscles}

  \begin{motsclesalter}
    \mot{users guide}
    \mot{J-SFdS document class}
  \end{motsclesalter}
  
  \begin{AMSclass}
    \kwd{35L05}
    \kwd{35L70}
  \end{AMSclass}
 
\end{frontmatter}
\section{Introduction}
Cette courte note décrit la classe de document \LaTeX \texttt{jsfds} pour préparer un manuscript à soumettre au Journal de la Société Française de Statistique. En lui-même c'est un exemple d'utilisation de cette classe. Il suppose que les auteurs ont une expérience de base avec \LaTeX ; dans le cas contraire, ils sont invités à se référer à \cite{lam94}.

\section{A propos du préambule et de la première page}
Votre fichier \LaTeX est essentiellement composé de deux parties : le préambule dans lequel vous insérez vos commandes personnelles se situe entre les commandes \verb+\documentclass+ et \verb+\begin{document}+ et le corps contenant le texte de l'article. Le corps du fichier est lui-même décomposé en deux parties : le descriptif de la première page (titre, auteurs, résumés, mots clefs) suivi du texte à proprement parlé. 

Le fichier \verb+jsfds-template-fr.tex+ contient une description de tous les champs que vous devez remplir afin d'obtenir un résultat satisfaisant.

\section{La structure du fichier \LaTeX}
\subsection{Préambule}
Il débute par les déclaration habituelles :
\begin{verbatim}
\documentclass[soumission]{jsfds}
\usepackage[latin1]{inputenc}
\usepackage[T1]{fontenc}
\end{verbatim}
dans lesquelles vous pouvez spécifier des encodages spécifiques de caractères : \verb+latin1+ ou un autre de la liste (\verb+ansinew applemac ascii latin9 utf8+)

Un champ est réservé pour introduire vos propres commandes \LaTeX :
\begin{verbatim}
\startlocaldefs
% par exemple
\newcommand{\R}{\mathbb{R}}
\newcommand{\E}{\mbox{\tiny \`eme}}
\def\1{1\kern-.20em {\rm l}}
\endlocaldefs
\end{verbatim}
Une commande déclare le langage principal de l'article, français ou anglais.
\begin{verbatim}
\setmainlanguagefrench
\end{verbatim}

\subsection{Première page}
Après ces deux commandes
\begin{verbatim}
\begin{document}
\begin{frontmatter}
\end{verbatim}
vous pouvez introduire toutes les informations requises pour décrire la première page de votre article :
\begin{verbatim}
\arxiv{math.PR/00000000}+ %si disponible 
\titre{Le titre complet}
\titrecourant{Titre résumé}
\titrealter{Traduction du titre en anglais}
\end{verbatim}
suivis des auteurs et de leur appartenance avec autant de commandes \verb+\auteur+ que d'auteurs.
\begin{verbatim}
\begin{InfosAuteurs}
    \auteur{%
      \prenom{Auteur} 
      \nom{Un}%
      \thanksref{t1}%
      \contact[label=e1]{auteur.un@math.univ.fr}%
    }%
    \and%
    \auteur{%
      \prenom{Auteur} 
      \nom{Deux}%
      \thanksref{t2}%
      \contact[label=e2]{auteur.deux@stat.edu}%
    }%
    \and%
    \auteur{%
      \prenom{Auteur} 
      \nom{Trois}%
      \thanksref{t2}%
      \contact[label=e3]{auteur.trois@stat.edu}%
    }
    \affiliation[t1]{Nom de l'institution auteur 1.\\ 
      \printcontact{e1}%
    }
    \affiliation[t2]{Nom de l'institution des auteurs 2 et 3.\\ 
      \printcontact{e2} and \printcontact{e3}% 
    }
    \enteteauteurs{Un, Deux et Trois} %noms courants des auteurs en tête de page
  \end{InfosAuteurs}
\end{verbatim}
Introduisez les résumés : 
\begin{verbatim}
\begin{abstract}
Résumé en français
\end{abstract}
\begin{altabstract}
Sa traduction en anglais
\end{altabstract}
\end{verbatim}
Les moths clefs sont rédigés en minuscules en français comme en anglais.
\begin{verbatim}
\begin{keywords}
\mot{premier}%
\mot{deuxième}%
\end{keywords}
\begin{altkeywords}
\mot{first}
\mot{second}
\end{altkeywords}
\end{verbatim}
Les sujets de classification avec la codification de l'AMS.
\begin{verbatim}
\begin{AMSclass}
\kwd{60K35}
\kwd{60L05}
\end{AMSclass}
\end{verbatim}
La commande suivante achève la première page. La taille des résumés est adaptée de façon à ne pas dépasser une page.
\begin{verbatim}
\end{frontmatter}
\end{verbatim}

\subsection{Corps de l'article}
Il est toujours possible de passer d'un texte en français (\verb+\selectlanguage{french}+) à un texte en anglais (\verb+\selectlanguage{english}+) tout en respectant les règles typographiques ainsi que les césures propres à chaque langage. 


\subsection{Inclure des graphiques}
Plusieurs librairies sont disponibles pour inclure des graphiques. Pour ce journal, vous êtes invité à utiliser la librairie \texttt{graphicx} développée par D.P. Carlisle et S.P.Q. Rahtz qui est déjà disponible dans votre installation de \LaTeX.
Suivez les instructions décrites en détail par \cite{gooms94} pour inclure des graphiques. Le Journal de la SFdS est en accès libre en ligne sous la forme de fichiers ``pdf''. Ces fichiers étant générés par la commande \texttt{pdflatex}, il est indispensable que les fichiers contenant des graphiques soient sous un format également ``pdf'' ou ``jpeg'' ; éviter les fichiers au format postscript. 

\subsection{Remerciements}
L'environnement \verb+\begin{acknowledgement} ... \end{acknowledgement}+ peut être utilisé immédiatement avant l'insertion de la bibliographie pour exprimer vos remerciements.
\begin{verbatim}
\begin{acknowledgement}
Le premier auteur est reconnaissant pour...
\end{acknowledgement}
\end{verbatim}
La commande(\verb+\thanks+) est générallement utilisée pour remercier des institutions tandis que l'environnement \texttt{acknowledgements} est réservé à des personnes.


\subsection{Références croisées et biblioraphie}
Dans tous les cas, les auteurs doivent utiliser les commandes \verb+\label, \ref, \pageref, \cite+. Chaque partie: section, équation, légende (commande \verb+\caption+)de tableau ou graphique doit être numérotée  et éventuellement référencée par une commande  \verb+\label{...}+, mais les équations non référencées dans le corps du texte n'ont pas à être numérotées.
Introduire des citations multiples sous la forme : \verb+\cite{A1,A2}+.

Il ya deux façons d'introduire la section des références bibliographiques. Soit par l'environnement  \verb+\thebibliography+ soit par un ficher au format \verb+\BibTeX+. Il est préférable d'utiliser \verb+\BibTeX+ comme ci-dessous mais en aucun cas une bibliographie construite ``à la main''.

\begin{verbatim}
\bibliography{biblio}
\end{verbatim}
Clore le fichier avec la commande \verb+\end{document}+.

\bibliography{biblio}
\end{document}